\usepackage[japanese]{babel}
\usepackage{luatexja}
\usepackage[ipa]{luatexja-preset}
\usepackage[a4paper,margin=2cm,truedimen]{geometry}
\usepackage{amsmath,amssymb}
\usepackage{booktabs}
\usepackage{caption}
\usepackage{listings}
\usepackage{xcolor}

\lstset{
    basicstyle=\ttfamily,
    keywordstyle=\color[RGB]{33,74,135}\bfseries,
    stringstyle=\color[RGB]{79,153,5},
    commentstyle=\color[RGB]{143,89,2}\itshape,
    %numberstyle=\footnotesize,
    %numbers=none,
    %stepnumber=1,
    %numbersep=15pt,
    backgroundcolor=\color[RGB]{240,240,240},
    frame=single,
    frameround=ffff,
    framesep=5pt,
    rulecolor=\color[RGB]{148,150,152},
    breaklines=true,
    breakautoindent=true,
    breakatwhitespace=true,
    breakindent=25pt,
    showspaces=false,
    showstringspaces=false,
    showtabs=false,
    tabsize=2,
    captionpos=t,
    linewidth=\textwidth,
    aboveskip=\medskipamount,
    belowskip=\medskipamount,
}

% 部（# → \section）の切れ目で改ページ（server-text の \pagebreak 運用に相当）
\usepackage{etoolbox}
\newcounter{cnpartbreak}
\setcounter{cnpartbreak}{0}
\pretocmd{\section}{%
  \ifnum\value{cnpartbreak}>0
    \clearpage
  \fi
  \stepcounter{cnpartbreak}%
}{}{}

    \usepackage{fancyhdr}
    \pagestyle{fancy}
    \fancyhead[L]{}
    \fancyhead[R]{\textsl{\rightmark}}
    \fancyfoot[LO,RE]{LPI-Japan}
    \fancyfoot[RO,LE]{https://linuc.org}
    \fancyfoot[CE,CO]{\thepage}

\renewcommand\contentsname{目次}


% 目次
\usepackage{tocloft}
\setlength{\cftsecnumwidth}{2em}
\setlength{\cftsubsecnumwidth}{3em}
\setlength{\cftsubsubsecnumwidth}{4em}
